\begin{pastexam}[label={px:ntu-110-comm-dsa-33}]{臺大電信所 110 學年度（B）第 33 題}
（單選；每題 2 分，答錯不倒扣）

Please choose the best possible answer in each of the following questions. In questions that ask for time complexities, please choose the tightest bound.

For problems 32 to 35, please consider a graph with $|V|$ vertices and $|E|$ edges, to which the Dijkstra algorithm is applied to find the shortest path. Assume $|E|$ is both $O(|V|^2)$ and $\Omega(|V|)$.

If Dijkstra algorithm is implemented with Fibonacci heap as priority queue, then the complexity is

A.~$O(|V|^2+|E|\log|V|)$ \qquad B.~$O(|V|^2)$ \qquad C.~$O(|E|\log|V|)$\\
D.~$O(|E|+|V|\log|V|)$ \qquad E.~$O(|E|)$
\end{pastexam}

\begin{pastsoln}
Dijkstra 演算法對優先佇列做 $|V|$ 次插入（或一開始建好）、$|V|$ 次 \proc{Extract-Min}、至多 $|E|$ 次
\proc{Decrease-Key}（每條邊鬆弛一次），總時間
\[
|V| \cdot T_{\text{extract}} + |E| \cdot T_{\text{decrease}} + (\text{建佇列})
\]
（資結書第~25.2~節）。

費氏堆積的 \proc{Extract-Min} 平攤 $O(\log|V|)$，\proc{Decrease-Key} 平攤 $O(1)$（資結書第~8.6~節），
總時間 $O(|V|\log|V| + |E|)$，選 D。 這是一連串操作的總時間，用平攤分析得到的界對整個演算法成立。

C 的 $O(|E|\log|V|)$ 是二元或二項堆積的時間，在 $|E| = \omega(|V|)$ 時比 D 大，不是最緊的。
E 的 $O(|E|)$ 太樂觀：$|E| = \Theta(|V|)$ 時它是 $O(|V|)$，但 Dijkstra 以比較為基礎依距離順序取出頂點，
需要 $\Omega(|V|\log|V|)$。
\end{pastsoln}
